\ifdefined\gkver\else\def\gkver{student}\fi
\documentclass[\gkver]{gaokaozhenti}
\begin{document}

\chapter{2002年上海理综}
%% paperType: 综合能力测试物理部分

\begin{enumerate}
\item
%% number: 10
%% paperName: 2002年上海理综第 10 题 3 分
%% typeId: 1
%% type: 单选题
%% chapter: 光学
%% point: 光的电磁说
%% method: 无
%% score: 3
%% degree: 750
%% duplicateId: 0
%% body:
下列有关各种家用电器性能的判断中，正确的是\xzanswer{C}

\fourchoices[answer=C]
{空调机和电风扇都可以降低室温}
{微波炉产生的微波是波长很短的机械波}
{电视机和收音机都能接收电磁波}
{电冰箱的耗电量夏天与冬天一样多}

%% answer: C
%% memo:
\memoanswer{
本题原卷及材料中未提供详解，待补充。
}

\item
%% number: 11
%% paperName: 2002年上海理综第 11 题 3 分
%% typeId: 1
%% type: 单选题
%% chapter: 电路
%% point: 电功电功率
%% method: 无
%% score: 3
%% degree: 650
%% duplicateId: 0
%% body:
下面列出了不同品牌的电视机、电风扇、空调机和电冰箱铭牌上的主要项目，试判断正常工作时，其中功率最大的是\xzanswer{D}

\begin{table}[htbp]
\centering
\footnotesize
\begin{tblr}{|c|c|c|c|}
\hline
$54\Ucm$ 彩色电视接收机 & BC-65B 电冰箱 & FS-69 电风扇 & KFR-33GW 空调机 \\
\hline
工作电压 $170\UV\sim240\UV$ & 额定电压 $220\UV$ & 规格 $400\Umm$ & 额定电压 $220\UV$ \\
\hline
工作频率 $50\UHz$ & 工作频率 $50\UHz$ & 额定电压 $220\UV$ & 工作频率 $50\UHz$ \\
\hline
额定功率 $85\UW$ & 额定功率 $70\UW$ & 工作频率 $50\UHz$ & 制冷/制热电流 $6.2\UA$/$6.2\UA$ \\
\hline
 & 耗电量 $0.50\UkWh$/$24\Uh$ & 额定功率 $65\UW$ & \\
\hline
\end{tblr}
\end{table}

\fourchoices[answer=D]
{电视机}
{电冰箱}
{电风扇}
{空调机}

%% answer: D
%% memo:
\memoanswer{
本题原卷及材料中未提供详解，待补充。
}

\item
%% number: 12
%% paperName: 2002年上海理综第 12 题 3 分
%% typeId: 1
%% type: 单选题
%% chapter: 电路
%% point: 电功电功率
%% method: 无
%% score: 3
%% degree: 650
%% duplicateId: 0
%% body:
根据上题铭牌上提供的信息，判断在12小时内正常使用的电冰箱与连续运转的电风扇消耗电能的情况是\xzanswer{B}

\fourchoices[answer=B]
{电冰箱比电风扇多}
{电风扇比电冰箱多}
{电风扇与电冰箱一样多}
{电冰箱比电风扇可能多，也可能少}

%% answer: B
%% memo:
\memoanswer{
本题原卷及材料中未提供详解，待补充。
}

\item
%% number: 13
%% paperName: 2002年上海理综第 13 题 3 分
%% typeId: 1
%% type: 单选题
%% chapter: 电路
%% point: 多用表的使用
%% method: 无
%% score: 3
%% degree: 750
%% duplicateId: 0
%% body:
如果收音机不能正常工作，需要判断干电池是否已经报废，可取出一节干电池用多用表来测量它的电动势，下列步骤中正确的是\xzanswer{BD}

①把多用表的选择开关置于交流 $500\UV$ 档或置于 OFF 档

②把多用表的红表笔和干电池的负极接触，黑表笔与正极接触

③把多用表的红表笔和干电池的正极接触，黑表笔与负极接触

④在表盘上读出电压值

⑤把多用表的选择开关置于直流 $25\UV$ 档

⑥把多用表的选择开关置于直流 $5\UV$ 档

\fourchoices[answer=BD]
{⑤③④①}
{②⑤①④}
{⑥③④①}
{⑥②④①}

%% answer: BD
%% memo:
\memoanswer{
本题原卷及材料中未提供详解，待补充。
}

\item
%% number: 14
%% paperName: 2002年上海理综第 14 题 3 分
%% typeId: 1
%% type: 单选题
%% chapter: 曲线运动
%% point: 圆周运动的应用
%% method: 无
%% score: 3
%% degree: 650
%% duplicateId: 0
%% body:
右图是上海锦江乐园新建的“摩天转轮”，它的直径达 $98\Um$；世界排名第五。游人乘坐时，转轮始终不停地匀速转动，每转一周用时 25 分钟，每个箱轿共有 6 个座位。试判断下列说法中正确的是\xzanswer{A}

\onepicture[w=4.5cm]{figs/14.jpg}

\fourchoices[answer=A]
{每时每刻，每个人受到的合力都不等于零}
{每个乘客都在做加速度为零的匀速运动}
{乘客在乘坐过程中对座位的压力始终不变}
{乘客在乘坐过程中的机械能始终保持不变}

%% answer: A
%% memo:
\memoanswer{
本题原卷及材料中未提供详解，待补充。
}

\item
%% number: 15
%% paperName: 2002年上海理综第 15 题 3 分
%% typeId: 1
%% type: 单选题
%% chapter: 曲线运动
%% point: 线速度角速度
%% method: 近似与估算
%% score: 3
%% degree: 750
%% duplicateId: 0
%% body:
右图是上海锦江乐园新建的“摩天转轮”，它的直径达 $98\Um$；世界排名第五。游人乘坐时，转轮始终不停地匀速转动，每转一周用时 25 分钟，每个箱轿共有 6 个座位。从图上可以估算出“摩天转轮”座位总数为\xzanswer{C}

\onepicture[w=4.5cm]{figs/15.jpg}

\fourchoices[answer=C]
{300座}
{336座}
{378座}
{408座}

%% answer: C
%% memo:
\memoanswer{
本题原卷及材料中未提供详解，待补充。
}

\item
%% number: 30
%% paperName: 2002年上海理综第 30 题 3 分
%% typeId: 3
%% type: 填空题
%% chapter: 功和能
%% point: 动能定理
%% method: 无
%% score: 3
%% degree: 750
%% duplicateId: 0
%% body:
足球守门员在发球门球时，将一个静止的质量为 $0.4\Ukg$ 的足球，以 $10\Ums$ 的速度踢出，这时足球获得的动能是 \tkanswer[1.5cm]{$20$}$\UJ$。足球沿草地作直线运动，受到的阻力是足球重力的 $0.2$ 倍，当足球运动到距发球点 $20\Um$ 的后卫队员处时，速度为 \tkanswer[1.5cm]{$4.5$（或 $\sqrt{20}$）}$\Ums$。（$g$ 取 $10\Umsq$）

\onepicture[w=6.0cm, align=right]{figs/30a.jpg}

%% answer: 
\jdanswer{足球获得的动能 $E_{\text{k}}=\frac{1}{2}mv^{2}=20\UJ$；到后卫队员处时的速度 $v'\approx4.5\Ums$（或 $\sqrt{20}\Ums$）。}

%% memo:
\memoanswer{
本题原卷及材料中未提供详解，待补充。
}

\item
%% number: 30
%% paperName: 2002年上海理综第 30 题 3 分
%% typeId: 1
%% type: 单选题
%% chapter: 力物体的平衡
%% point: 受力分析
%% method: 无
%% score: 3
%% degree: 850
%% duplicateId: 0
%% body:
足球运动员已将足球踢向空中，如右图。下列描述足球在向斜上方飞行过程中某时刻的受力图中，正确的是\xzanswer{B}

\fourchoices[answer=B, ispicture=true, h=2.0cm]
{figs/30b_A.png}
{figs/30b_B.png}
{figs/30b_C.png}
{figs/30b_D.png}

（$G$ 为重力，$F$ 为脚对球的作用力，$F_{\text{阻}}$ 为空气阻力）

%% answer: B
%% memo:
\memoanswer{
本题原卷及材料中未提供详解，待补充。
}

\item
%% number: 38
%% paperName: 2002年上海理综第 38 题 4 分
%% typeId: 3
%% type: 填空题
%% chapter: 直线运动
%% point: 匀速直线运动
%% method: 无
%% score: 4
%% degree: 650
%% duplicateId: 0
%% body:
节水灌溉的方式有喷灌、滴灌等。滴灌是通过不同口径的塑料管，将水直接送到每株农作物的根部，以点滴等方式进行灌溉。现有一块农田共装了600个滴水头，如图所示，总进水管为 $p$，每个滴水头为 $q$。设总进水管横截面积 $S$ 为 $0.0016\Umq$，平均每分钟每个滴头滴水120滴（每5滴水体积约为 $1\Ucmc$），则总水管 $p$ 中的水流量是 \tkanswer[2cm]{$2.4\times10^{-4}$}$\Umcs$，$p$ 管中水流平均速度是 \tkanswer[2cm]{$0.15$}$\Ums$（流体在每秒钟流过某横截面的体积叫流量 $Q$，它与水管的横截面积 $S$ 及水流平均速度 $v$ 的关系是 $Q=Sv$）。

\onepicture[w=7.0cm, align=right]{figs/38.jpg}

%% answer: 
\jdanswer{总水管 $p$ 中的水流量 $Q=2.4\times10^{-4}\Umcs$；$p$ 管中水流平均速度 $v=0.15\Ums$。}

%% memo:
\memoanswer{
本题原卷及材料中未提供详解，待补充。
}

\item
%% number: 41
%% paperName: 2002年上海理综第 41 题 6 分
%% typeId: 3
%% type: 填空题
%% chapter: 气体性质
%% point: 能源的开发与利用
%% method: 无
%% score: 6
%% degree: 550
%% duplicateId: 0
%% body:
“东方绿舟”内有一个绿色能源区，同学们可以在这里做太阳能和风能的研究性实验。某同学为了测定夏季中午单位面积上、单位时间内获得的太阳能，制作了一个太阳能集热装置，实验器材有：（A）内壁涂黑的泡沫塑料箱一个，底面积为 $1\Umq$；（B）盛水塑料袋一个；（C）温度计一个；（D）玻璃板一块（约 $1\Umq$）。如图（a）。

\onepicture[w=13.0cm]{figs/41.png}

\begin{enumerate}
	\item
	假设图（b）为一斜坡草地，太阳光垂直照射到草地表面，请将上述实验器材按实验设计要求画在图（b）中。
	\item
	如果已知水的比热容为 $c$，被水吸收的热量 $Q$ 与水的质量 $m$、水温升高量 $\Delta t$ 间的关系是 $Q=cm\Delta t$，则为了测定中午单位面积上、单位时间内获得的太阳能，除了需要测量 $m$、$\Delta t$ 外，还应测量的物理量是 \tkanswer[2cm]{时间 $T$（或 $t$）}。本实验会有一定误差，试写出一条产生误差的主要原因：\tkanswer[6cm]{太阳能没有全部被水吸收，或水吸收的太阳能还有一部分损失等}。
\end{enumerate}

%% answer: 
\jdanswer{
\begin{enumerate}
	\item 
	实验设计装配图如图所示。
	
	\item
	时间 $T$（或 $t$）；太阳能没有全部被水吸收，或水吸收的太阳能还有一部分损失等。
\end{enumerate}
}

%% memo:
\memoanswer{
（1）实验设计装配图如图所示（将涂黑的泡沫塑料箱置于斜坡草地上，箱内放盛水塑料袋，盖上玻璃板，使太阳光垂直照射到草地表面；两种画法均可）。

\onepicture[w=6.5cm]{figs/41_an.jpg}

（2）由 $Q=cm\Delta t$ 可求出水吸收的热量，测出受热时间 $T$，再由集热面积 $S$ 即可求出单位面积上、单位时间内获得的太阳能，故还应测量的物理量是时间 $T$（或 $t$）。

（3）产生误差的主要原因：太阳能没有全部被水吸收，或水吸收的太阳能还有一部分损失等。
}

\item
%% number: 46
%% paperName: 2002年上海理综第 46 题 4 分
%% typeId: 5
%% type: 辨析题
%% chapter: 运动和力
%% point: 牛顿第二定律
%% method: 无
%% score: 4
%% degree: 650
%% duplicateId: 0
%% body:
在水平路面上，一个大人推一辆重车，一个小孩推一辆轻车，各自作匀加速直线运动（阻力不计），甲、乙两同学在一起议论。甲同学说：根据牛顿运动定律，大人的推力大，小孩的推力小，因此重车的加速度大。乙同学说：根据牛顿运动定律，重车质量大，轻车质量小，因此轻车的加速度大。

你认为他们的说法是否正确？请简述理由。

%% answer: 
\jdanswer{两种说法都不对。理由见详解。}

%% memo:
\memoanswer{
这两种说法都不对。根据牛顿第二定律 $a=\frac{F}{m}$，只有 $m$ 一定时，才能根据 $F$ 的大小来确定 $a$ 的大小，同样只有 $F$ 一定时才能根据 $m$ 的大小确定 $a$ 的大小，在两者都不相同的情况下，只能用 $F$ 与 $m$ 的比值来确定 $a$ 的大小。
}

\item
%% number: 49
%% paperName: 2002年上海理综第 49 题 6 分
%% typeId: 3
%% type: 填空题
%% chapter: 电路
%% point: 闭合电路欧姆定律
%% method: 无
%% score: 6
%% degree: 550
%% duplicateId: 0
%% body:
研究性学习是一种重要的学习方式，有利于培养学生的实践能力和创新精神。

“M3 型电池和普通碱性电池性价比的研究”是某校学生摄影协会提出的一个研究课题，目的是通过研究，为摄影协会选择性价比（可拍摄照片张数与所用电池费用之比）较高的电池提供依据。他们分成 A、B 两组开展了调查和实验，A 小组利用两台相同型号的照相机，分别装上两种电池各 4 节在暗室中连续拍照，每次都用闪光灯，直到两种电池都不能使相机闪光为止。B 小组用同样装置，在校外不同地方拍摄，有时用闪光灯，有时不用。结果 A 小组用 M3 电池拍了 290 张照片；用普通碱性电池拍了 184 张照片。B 小组用 M3 型电池拍了 583 张照片；用普通碱性电池拍了 364 张照片。

你认为 \tkanswer[1cm]{A} 小组实验结果更可靠，原因是他们控制了 \tkanswer[4cm]{每次拍照都闪光和拍照时背景亮度} 等影响实验结果的因素。如果 M3 型电池每4节36元，普通碱性电池每4节28元，试参照左下直方图，画出该组的性价比直方图。

\onepicture[w=13.0cm]{figs/49.png}

你认为他们的实验还存在的问题是（举出一例）\tkanswer[4cm]{实验所选用的电池样本太少}。

请你提出一个与 M3 型电池有关的新的、有价值的课题名称：\tkanswer[4cm]{M3电池工作寿命较长的原因探究}。

%% answer: 
\jdanswer{
\begin{enumerate}
	\item
	A；每次拍照都闪光和拍照时背景亮度。
	
	\item
	实验所选用的电池样本太少。如“M3电池工作寿命较长的原因探究”“M3电池的放电规律实验研究”“M3电池的环保问题研究”等等。
\end{enumerate}
}

%% memo:
\memoanswer{
（1）A 小组实验结果更可靠，因为他们控制了每次拍照都闪光和拍照时背景亮度等影响实验结果的因素。

A 组 M3 型电池的性价比为 $\frac{290}{36}\approx8.1$（张/元），普通碱性电池的性价比为 $\frac{184}{28}\approx6.6$（张/元），该组的性价比直方图如图所示。

\onepicture[w=6.0cm]{figs/49_an_A.png}

（2）实验所选用的电池样本太少。可提出的新课题如“M3电池工作寿命较长的原因探究”“M3电池的放电规律实验研究”“M3电池的环保问题研究”等等。

（说明：原卷“A 小组用 M3 电池拍了 2 对张照片”中“2对张”系原卷排版缺字，据左下拍摄张数直方图应为 290 张。）
}

\end{enumerate}

\end{document}
